\chapter{Introduction}
\label{sec:intro}

Reverberation mapping (RM) of active galactic nuclei measures the delay between continuum and broad-line variations
and therefore needs spectra whose wavelength scale, line-spread function and relative flux calibration are
consistent from epoch to epoch at the per-cent level. The narrow emission lines (in particular [O\,{\sc iii}]
$\lambda\lambda4959,5007$) are used afterwards to put all epochs on a common flux scale and to align them in
wavelength \citep{vanGroningen1992,Fausnaugh2017}, but this inter-calibration can only remove grey
offsets, small shifts and small resolution differences: wavelength-dependent errors on scales of tens to hundreds of
\AA{} (fringing, sensitivity-function structure, order mismatches, extrapolated wavelength solutions) propagate
directly into the broad-line light curves.

The FLOYDS spectrographs at Las Cumbres Observatory (LCO) record two cross-dispersed orders simultaneously: the red
order (order~1, $\approx4500$--$10\,000$\,\AA, $\approx3.5$\,\AA\,pixel$^{-1}$) and the blue order (order~2,
$\approx3200$--$5800$\,\AA, $\approx1.7$\,\AA\,pixel$^{-1}$), with a spatial scale of 0.337\,arcsec\,pixel$^{-1}$.
RM campaigns use the 6\,arcsec slit to minimise slit losses. The archive products are reduced by the
\banzai{} pipeline. Members of our collaboration found the wavelength calibration, the extraction and the flux
calibration of these products unsatisfactory for RM. This report documents a systematic diagnosis of the archive
products of \banzai{}~1.0.1 (the version that produced all data used here), traces each problem to its origin in the
code or the calibration data, and describes target-independent corrections, now built into \banzai{} and run from the
raw frames (improved v3, Table~\ref{tab:versions_names}; the names of the reductions and other recurring terms are
defined in the Conventions, p.~\pageref{sec:conventions}). They were first developed on the public
\texttt{-1d}/\texttt{-2d} archive products, which required undoing parts of the 1.0.1 processing; that procedure
(improved v1) is documented in Appendix~\ref{sec:app_v1}. All diagnostics and corrections are implemented as reusable
modules inside the \banzai{} package (Chapter~\ref{sec:software}).

Throughout, ``target-independent'' means that no information about the science target (redshift, line list, spectral
shape) is used to derive a correction. The known redshift of Fairall~9 and its narrow lines are used only as an
external check.

\section{Data}
\label{sec:data}

\begin{table}
\centering
\caption{Data used in this report. All frames are from FLOYDS on the 2-m Faulkes Telescope South at Siding Spring
(LCO telescope code \texttt{coj2m002}, camera \texttt{en12}), reduced by \banzai{}~1.0.1.}
\label{tab:data}
\begin{tabular}{lccc}
\toprule
Set & $N$ & Slit & Dates (2026) \\
\midrule
Fairall~9 (proposal KEY2026B-007), 600\,s & 50 & 6\,arcsec & Sep 25 -- Nov 11 \\
EG\,21 (``FLOYDS standards''), 300\,s & 18 & 2\,arcsec & Oct 10 -- Nov 12 \\
Raw arcs (Hg + Ar lamps; Zn listed as on), 60\,s & 64 & 6\,arcsec & Aug 1 -- Nov 11 \\
\bottomrule
\end{tabular}
\end{table}

Table~\ref{tab:data} lists the data. For every science and standard exposure we use the archive \texttt{-1d}
product (FITS extensions \texttt{SPECTRUM}, \texttt{EXTRACTED}, \texttt{SENSITIVITY}, \texttt{TELLURIC}) and the
\texttt{-2d} product (image, error, bad-pixel mask, order map, wavelength map, fringe map, object profile and sky
model: \texttt{SCI}, \texttt{ERR}, \texttt{BPM}, \texttt{ORDERS}, \texttt{WAVELENGTH}, \texttt{FRINGE},
\texttt{PROFILE}, \texttt{BACKGROUND}; and the pixel table \texttt{BINNED2D}). The pixel table lists
every pixel of both orders with its wavelength, its position along the slit relative to the order centre
($y_{\rm o}$) and to the fitted trace ($y_{\rm p}$), the profile width $\sigma_{\rm p}$, the column-normalised
profile weight $P$ and the 1-pixel wavelength bin it belongs to. We verified that \texttt{BINNED2D} reproduces
\texttt{SCI}, \texttt{ERR}, \texttt{BACKGROUND} and \texttt{PROFILE} to the quantisation of the compressed images.
For each science frame, the header keyword \texttt{L1IDARC} names the processed arc (file type \texttt{a91}) that the
pipeline used; our raw arcs are the corresponding raw files (\texttt{a00}). For 49 of the 50 science frames that arc
is in our set. The arcs of the standards (taken by another
programme) were not available.

Fairall~9 was observed at airmass 1.13--1.85 (median 1.21; 34 of 50 frames below 1.3). We adopt
$z=0.047016$ for the external checks.

These archive products are used for the diagnosis (Chapters~\ref{sec:traces}--\ref{sec:flux}). The reduction from
raw uses the full COJ campaign (412 raw frames, Chapter~\ref{sec:raw}), and the final version is verified on a second
campaign, NGC\,4151 with the FLOYDS on the Faulkes Telescope North (OGG; 421 raw frames, Chapter~\ref{sec:sites}).
